\subsection{A conta de três participantes}
\label{subsec:txtangle-tres-participantes}

Consideremos Alice, Bob e Clara, designados por \(A,B,C\). Cada um chega
com entradas, destinos e uma parte acordada da taxa. Para seguir a
aritmética dos escalares, façamos primeiro as contas módulo 19, como se
o grupo tivesse essa ordem prima. Os três ajustes escolhidos pelos pares
são
\[
              s_{A,B}=4,\qquad s_{A,C}=9,\qquad s_{B,C}=3 .
\]
A convenção de sinais dá ao primeiro membro do par o valor escrito
e ao segundo o seu simétrico. A soma por participante lê-se melhor
quando a dispomos por colunas.
\begin{center}
\begin{tabular}{c|rrr|r}
 &par \(A,B\)&par \(A,C\)&par \(B,C\)&soma módulo 19\\
\hline
Alice &\(4\)&\(9\)&\(0\)&\(13\)\\
Bob &\(-4\)&\(0\)&\(3\)&\(18\)\\
Clara &\(0\)&\(-9\)&\(-3\)&\(7\)
\end{tabular}
\end{center}

Bob tem \(-1\), cujo representante canónico é 18, e Clara tem \(-12\),
cujo representante é 7. Todos os termos de uma coluna somam zero. Nas
linhas, a verificação dos representantes canónicos exige uma redução:
\(13+18+7=38\equiv0\pmod{19}\). É lícito escolher qualquer
representante quando estes números multiplicam \(G\), pois
\([19]G=\mathcal O\) no grupo do exemplo.

Suponhamos que as três diferenças locais entre pseudo-compromissos de
entrada e compromissos de saída, já descontada a fracção de taxa de
cada um, davam \(\mathcal O\) antes do ajuste. Depois da adição dos
três termos passam a dar \(13G,18G,7G\), respectivamente. Pergunte-se
agora quanto vale a soma de Alice e Bob. Pelas linhas vale
\((13+18)G=12G\). Pelas colunas, o acordo entre eles entra uma vez como
\(4G\) e outra como \(-4G\), e desaparece. Sobram apenas os dois
acordos com Clara, \((9+3)G=12G\). Ao juntar a parte de Clara
obtemos \((12+7)G=\mathcal O\). Vê-se no mesmo cálculo como a soma
parcial conserva o contributo dos pares que atravessam a fronteira
escolhida e como a soma total o perde.

Há uma diferença entre fazer esta conta e saber antecipadamente quanto
ela dará. Para a perceber, conservemos o 4 do par \(A,B\), mas
representemos os números dos outros dois pares por \(u,v\). Os ajustes
passam a ser
\[
            (\delta_A,\delta_B,\delta_C)
                  =(4+u,\,-4+v,\,-u-v).
\]
Se \(u,v\) forem independentes e uniformes módulo 19, os dois
primeiros ajustes percorrem os \(19^2=361\) pares possíveis com
igual probabilidade. Para obter um par determinado \((a,b)\) há
exactamente uma escolha, \(u=a-4\) e \(v=b+4\). A terceira coordenada
fica então obrigada a \(-a-b\). Dos \(19^3\) triplos de escalares,
precisamente 361 satisfazem a soma nula, e todos esses triplos são
igualmente prováveis neste modelo.

A soma de Alice e Bob vale agora \((u+v)G\). Mesmo que \(u\) esteja
fixado, a parcela continuará uniforme se \(v\) for uniforme perante a
informação de quem observa. Para cada alvo \(wG\) há exactamente um
\(v=w-u\). A probabilidade de atingir esse alvo é \(1/19\) e,
em particular, a soma parcial pode dar a identidade. A soma dos três,
essa sim, dá sempre a identidade. O número 19 deixa contar todas
as escolhas. No grupo do protocolo o mesmo argumento substitui 19
pela ordem \(l\), conservando a hipótese sobre o que permanece
desconhecido a quem faz a observação.

\paragraph{Quem paga o quê.}
A conta em \(G\) resolveu a distribuição das máscaras. Os montantes
têm uma contabilidade própria, no coeficiente de \(H\). Fixemos,
em unidades inteiras, a seguinte proposta. Alice contribui com uma
entrada de 11, envia 7 e recupera 3. Bob contribui com 8, envia 5
e recupera 2. Clara contribui com 6, envia 4 e recupera 1. Cada
um suporta uma unidade da taxa. Por participante temos
\[
                 11=7+3+1,\qquad 8=5+2+1,\qquad 6=4+1+1 .
\]
No conjunto, entram 25, formam-se saídas de \(10+7+5=22\) e
declara-se uma taxa de 3. Estas quantidades dizem respeito ao grupo
de ordem \(l\) da transacção e aos intervalos inteiros admissíveis
para os montantes. O pequeno módulo anterior servia apenas para pôr
à vista a soma das máscaras.

Se Alice deslocar o seu pseudo-compromisso em \(\delta_A G\), o
coeficiente de \(H\) da entrada conserva o valor 11. Por isso pode
conferir que o destinatário recebe 7 e que lhe cabem 3 de troco antes
de autorizar a entrada. Bob e Clara fazem a mesma conferência sobre as
suas saídas. O nó recebe a transacção completa, com a taxa global e
um balanço de pontos. Já a distribuição das três unidades de taxa por
quem colaborou pertence ao acordo dos participantes. A distinção entre
essas duas vistas determina que informação deve circular nas rondas
de construção que vamos acompanhar.
